\section{Background}
\label{sec:preliminaries}

We set $\N \coloneqq \set{1,2,\dots}$ and always assume that $n\in \N$. For $a,b\in \Z$ we define the \emph{discrete interval} $[a,b]\coloneqq \set{c\in \Z\mid a\leq c \leq b}$ and may use the shorthand $[a] \coloneqq [1,a]$. 
Throughout this paper~$\K$ denotes an arbitrary field. For a set $X$, $\spa_\K X$ is the $\K$-vector space with basis~$X$.

\subsection{Symmetric groups and 0-Hecke algebras}

The \emph{symmetric group} $\fS_n$ is the group of all permutations of the set~$[n]$. We proceed by reviewing $\fS_n$ as a Coxeter group. More details can be found in~\cite{Bjorner2006}.

As a Coxeter group $\fS_n$ is generated by the adjacent transpositions $\fs_i \coloneqq (i,i+1)$ for $i= 1,\dots, n-1$ which satisfy
\begin{alignat*}{3}
\fs_i^2 &= 1, 			\\
\fs_i\fs_{i+1}\fs_i &= \fs_{i+1}\fs_i\fs_{i+1}, 	\\
\fs_i\fs_j &= \fs_j \fs_i 		\ \text{if } |i-j|\geq 2. 
\end{alignat*}
The latter two relations are called \emph{braid relations}.
Let $\sigma \in \fS_n$. We can write $\sigma$ as a product $\sigma = s_{j_k} \cdots s_{j_1}$.
If~$k$ is minimal among such expressions, $s_{j_k} \cdots s_{j_1}$ is a \emph{reduced word} for $\sigma$ and $\ell(\sigma) \coloneqq k$ is the \emph{length} of $\sigma$. 

The \emph{support} of $\sigma$ is $\supp(\sigma) =\{ i\in [n-1] \mid \text{$\fs_i$ appears in a reduced word of $\sigma$}\}$.
One assertion of the \emph{word property} of $\fS_n$~\cite[Theorem~3.3.1]{Bjorner2006} is that a reduced word of $\sigma$ can be transformed into any other reduced word of $\sigma$ by applying a sequence of braid relations. 
Thus, for each reduced word of $\sigma$ the set of indices occurring in it is $\supp(\sigma)$.

Let $\sigma,\tau \in \fS_n$. The \emph{left weak order} $\leq_L$ is the partial order on $\fS_n$ given by
\begin{equation*}
\sigma \leq_L \tau \iff 
\begin{aligned}
&\tau = \fs_{i_k} \cdots \fs_{i_1} \sigma, \\
&\ell(\fs_{i_r} \cdots \fs_{i_1} \sigma) = \ell(\sigma) + r \text{ for } r=1,\dots ,k.
\end{aligned}
\end{equation*}
The following~\cref{thm:properties_bruhat_order} gathers immediate consequences of the definition.

\begin{prop} 
		\label{thm:properties_bruhat_order}
	Let $\sigma,\tau \in \fS_n$.
	\begin{enumerate}
		\item We have $\sigma \leq_L \tau$ if and only if $\ell(\tau\sigma^{-1}) = \ell(\tau) - \ell(\sigma)$.
		\item If $\sigma\leq_L \tau$ then the reduced words for $\tau\sigma^{-1}$ are in bijection with saturated chains in the left weak order poset $(\fS_n, \leq_L)$ from $\sigma$ to $\tau$ via 
		\begin{align*}
		\fs_{i_k} \cdots \fs_{i_1} \quad \longleftrightarrow \quad \sigma \lessdot_L \fs_{i_1} \sigma \lessdot_L \fs_{i_2}\fs_{i_1} \sigma \lessdot_L\dots \lessdot_L \fs_{i_k} \cdots \fs_{i_1}\sigma = \tau.
		\end{align*}
		\item The poset $(\fS_n, \leq_L)$ is graded by the length function.
	\end{enumerate}
\end{prop}

\begin{thm}[{\cite[Corollary~3.2.2]{Bjorner2006}}]
	\label{thm:bruhat_interval}
	Let $\sigma,\tau\in \fS_n$. The interval in left weak order $[\sigma, \tau] \coloneqq \left\{\rho \in \fS_n \mid \sigma \leq_L \rho \leq_L \tau \right\}$ is a graded lattice with rank function $\rho \mapsto \ell(\rho \sigma^{-1})$.
\end{thm}
 
Next, we define the 0-Hecke algebra~$\He$. We use the presentation as in~\cite{Tewari2015} and refer to~\cite[Ch.~1]{Mathas1999} for details.

\begin{defi}
	\label{thm:0-Hecke algebra}
	The \emph{0-Hecke algebra}~$\He$ is the unital associative $\K$-algebra generated by the elements $\pi_1,\pi_2,\dots ,\pi_{n-1}$ subject to relations
	\begin{align*}
	\pi_i^2 &= \pi_i, \\
	\pi_i\pi_{i+1}\pi_i &= \pi_{i+1} \pi_{i} \pi_{i+1}, \\
	\pi_i \pi_j &= \pi_j \pi_i \text{ if } |i-j| \geq 2. 	 
	\end{align*}
\end{defi}

Note that the $\pi_1,\dots, \pi_{n-1}$ are projections satisfying the same braid relations as the $s_1, \ldots, s_{n-1}$.
Another set of generators of~$\He$ is given by elements $\bar \pi_i$ for $i=1,\dots, n-1$ such that $\bar \pi_i^2 = - \bar \pi_i$ and the braid relations hold.
We do not use them here but they appear in some of the references.
%\cite{Norton1979, Mathas1999, Huang2016, Duchamp1996, Krob1997}.

Let $\sigma \in \fS_n$. We define $\pi_\sigma \coloneqq \pi_{j_k} \cdots \pi_{j_1}$ where $\fs_{j_k} \cdots \fs_{j_1}$ is a reduced word for $\sigma$. The word property ensures that this is well defined. Multiplication is given by
\begin{align*}
\pi_i \pi_\sigma = 
\begin{cases}
\pi_{\fs_i \sigma} & \text{if}\ \ell(\fs_i \sigma ) > \ell(\sigma) \\
\pi_{\sigma} & \text{if}\ \ell(\fs_i \sigma ) < \ell(\sigma) \\
\end{cases}
\end{align*}
for $i = 1,\dots, n-1$. 
%\cite[Lemma 1.12]{Mathas1999}. 
As a consequence, $\set{\pi_{\tau} \mid \tau \in \fS_n}$ spans~$\He$ over~$\K$. One can also show that it is a $\K$-basis of~$\He$.%\cite[Theorem 1.13]{Mathas1999}.


\subsection{Compositions and composition tableaux}
\label{sec:compositions_and_composition_tableaux}

A \emph{composition} $\alpha = \parts{\alpha}l$ is a finite sequence of positive integers. 
The \emph{length} and the \emph{size} of $\alpha$ are given by $\ell(\alpha)\coloneqq l$ and $|\alpha| \coloneqq \sum_{i=1}^l \alpha_i$, respectively.
The $\alpha_i$'s are called \emph{parts} of $\alpha$. If $\alpha$ has size~$n$, $\alpha$ is called \emph{composition of~$n$} and we write $\alpha \vDash n$.
A \emph{partition} is a composition whose parts are weakly decreasing. 
We write $\lambda \vdash n$ if $\lambda$ is a partition of size~$n$.
For a composition $\alpha$ we denote the partition obtained by sorting the parts of $\alpha$ in decreasing order by $\widetilde \alpha$. The \emph{empty composition} $\emptyset$ is the unique composition of length and size~$0$.

\begin{exam} 
	For $\alpha=(1,4,3)\vDash 8$, we have $\widetilde \alpha = (4,3,1)\vdash 8$.
\end{exam}

A \emph{cell} $(i,j)$ is an element of $\N \times \N$.
A finite set of cells is called \emph{diagram}. Diagrams are visualized in English notation. That is, for each cell $(i,j)$ of a diagram we draw a box at position $(i,j)$ in matrix coordinates.
The \emph{diagram} of $\alpha\vDash n$ is the set $\set{(i,j)\in \N \times \N \mid i \leq \ell(\alpha), j \leq \alpha_i}$. So, we display the diagram of $\alpha$ by putting $\alpha_i$ boxes in row~$i$ where the top row has index~$1$.
We often identify $\alpha$ with its diagram.

\begin{exam}
	\begin{align*}
		(1,4,3) = 
		\begin{ytableau} 
		 \\
			&&& \\
			&&\\
		\end{ytableau}
	\end{align*}
\end{exam}

Next, we will consider standard composition tableaux and a related poset of compositions which arose in~\cite{Bessenrodt2011}.

\begin{defi}
	The \emph{composition poset} $\mc L_c$ is the set of all compositions together with the partial order $\leq_c$ given as the transitive closure of the following covering relation. For compositions $\alpha$ and $\beta = ( \beta_1, \ldots, \beta_l)$
	\begin{equation*}
		 \beta \lessdot_c \alpha \iff \begin{aligned}
		 \alpha &= (1, \beta_1, \ldots, \beta_l) \text{ or}\\
		 \alpha &= ( \beta_1, \ldots, \beta_k +1 , \ldots, \beta_l) \text{ and $ \beta_i \neq \beta_k$ for all $i<k$.}
		\end{aligned} 
	\end{equation*}	
\end{defi} 

In other words, $\beta$ is covered by $\alpha$ in $\mc L_c$ if and only if the diagram of $\alpha$ can be obtained from the diagram of $\beta$ by adding a box as the new first row or appending a box to a row which is the topmost row of its length in $\beta$.

\begin{exam}
	\label{exa:chain_in_composition_poset}
	\[	
		\ydiagram{1,2} \lessdot_c 
		\ydiagram{2,2} \lessdot_c 
		\ydiagram{3,2} \lessdot_c 
		\ydiagram{3,3} \lessdot_c 
		\ydiagram{1,3,3} \lessdot_c
		\ydiagram{1,4,3} 
	\]
\end{exam}

Let $\alpha$ and $\beta$ be two compositions such that $\beta \leq_c \alpha$. In this situation we always assume that the diagram of $\beta$ is moved to the bottom of the diagram of $\alpha$, and we define the \emph{skew composition diagram} (or \emph{skew shape})~$\ab$ to consist of all cells of $\alpha$ which are not contained in $\beta$.

Moreover, we define $\osh(\ab)= \alpha$ and $\ish(\ab)= \beta$ as the \emph{outer} and the \emph{inner shape} of~$\ab$, respectively. 
The size of a skew shape is $|\ab| \coloneqq |\alpha| - |\beta|$.
We call~$\ab$ \emph{straight} if $\beta = \emptyset$.
In this case the skew composition diagram~$\ab$ is nothing but the ordinary composition diagram $\alpha$.

\begin{exam} In the following the cells of the inner shape are gray.
	\[
		(1,4,3) \skc (1,2) =
		\begin{ytableau} 
			 \\
			*(gray)&		&& \\
			*(gray)&*(gray)	&\\
		\end{ytableau}
	\]	
\end{exam}

Note that $\beta \leq_c \alpha$ implies $\beta_{\ell(\beta)-i} \leq \alpha_{\ell(\alpha) - i}$ for $i=0,\dots, \ell(\beta) - 1$. One could define skew shapes for all pairs of compositions fulfilling this containment condition.
Anyway, we demand $\leq_c$ rather than containment since with the latter one allows skew shapes for which standard composition tableaux (which we will define next) do not exist.
For instance, the compositions $\beta = (1,1)$ and $\alpha = (1,2)$ satisfy the containment condition but $\beta \not \leq_c \alpha$.
Even if~$\ab$ were a skew shape, there would be no standard composition tableau of this shape (because of the triple rule stated below).

Let~$D$ be a diagram. A \emph{tableau}~$T$ of shape~$D$ is a map $T\colon D \to \N$. It is visualized by filling each $(i,j)\in D$ with $T(i,j)$.

\begin{defi}
	Let~$\ab$ be a skew shape of size~$n$. A \emph{standard composition tableau} (SCT) of shape~$\ab$ is a bijective filling $T\colon \ab \to [n]$ satisfying the following conditions:
	\begin{enumerate}
		\item
		The entries are decreasing in each row from left to right.
		\item
		The entries are increasing in the first column from top to bottom.		
		\item
		(Triple rule). Set $T(i,j)\coloneqq \infty$ for all $(i,j)\in \beta$.
		If $(j,k)\in \alpha \skc \beta$ and $(i,k-1)\in \alpha$ such that $j > i$ and $T(j,k) < T(i,k-1)$ then $(i,k) \in \alpha$ and $T(j,k) < T(i,k)$.

	\end{enumerate}
\end{defi}

Let $a \coloneqq T(j,k)$, $b\coloneqq T(i,k-1)$ be two entries of an SCT~$T$ occurring in adjacent columns. Then the triple rule can be visualized as follows by considering the positions of entries in~$T$:
\begin{align*}
	\begin{ytableau}
		b \\
		\none \\
		\none \\ 
		\none & a
	\end{ytableau}
	\quad
	\text{and $a<b$}
	\quad \overset{\textit{triple rule}}\implies \quad
	\text{$\exists c\in T \colon$ }
	\begin{ytableau}
		b & c \\
		\none \\
		\none \\ 
		\none & a
	\end{ytableau}
	\quad
	\text{and $a<c$}.
\end{align*}

The set of standard composition tableaux of shape~$\ab$ is denoted with $\SCT(\ab)$. 
For an SCT~$T$ we write $\sh(T)$ for its shape. 
The notions of outer and inner shape are carried over from $\sh(T)$ to~$T$.
We call~$T$ \emph{straight} if its shape is straight.

\begin{exam}
	\label{exa:sct}
	 An SCT is shown below.
	\begin{align*}
	T = 
	\begin{ytableau} 
		2 \\ 
		*(gray) & 5 & 4 & 1 \\ 
		*(gray) & *(gray) & 3 \\ 
		\end{ytableau}
	\end{align*}
	We have $\osh(T)= (1,4,3)$ and 	$\ish(T) = (1,2)$.
\end{exam}

Standard composition tableaux encode saturated chains of $\mathcal{L}_c$ in the following way.

\begin{prop}[see {\cite[Proposition~2.11]{Bessenrodt2011}}]
	\label{thm:sct_and_chains}
	Let~$\ab$ be a skew composition of size~$n$. For $T\in \SCT(\ab)$, 
	\begin{align*}
		\beta = \alpha^n \lessdot_c \alpha^{n-1} \lessdot_c \dots \lessdot_c \alpha^0 =\alpha
	\end{align*}
	given by
	\begin{align} \label{eq:tableau_chain}
	\alpha^n = \beta, \quad \alpha^{k-1} = \alpha^k \cup T^{-1}(k) \quad \text{for} \quad k= 1,\dots, n
	\end{align}
	is a saturated chain in $\mathcal{L}_c$. Moreover, we obtain a bijection from $\SCT(\ab)$ to the set of saturated chains in $\mathcal{L}_c$ from $\beta$ to $\alpha$ by mapping each tableau of $\SCT(\ab)$ to its corresponding chain given by~\eqref{eq:tableau_chain}. 	 
\end{prop}

\begin{exam}
	The SCT from~\cref{exa:sct} 	corresponds to the chain from~\cref{exa:chain_in_composition_poset}.	
\end{exam}

From the perspective of~\cref{thm:sct_and_chains}, the triple rule reflects the fact that by adding a cell to a row of a composition diagram, a covering relation in $\mc L_c$ is established if and only if the row in question is the topmost row of its length.

Some of the upcoming notions already played a role in~\cite{Tewari2015}.
Let $(i,j)$ and $(i',j')$ be two cells.
Define $r(i,j) \coloneqq i$ and $c(i,j) \coloneqq j$ the \emph{row} and the \emph{column} of $(i,j)$, respectively. 
We say that $(i,j)$ \emph{attacks} $(i',j')$ and write $(i,j)\att (i',j')$ if $j=j'$ and $i\neq i'$ or $j = j'-1$ and $i<i'$.
That is, the two cells are distinct and either they appear in the same column or they appear in adjacent columns such that $(i',j')$ is located strictly below and right of $(i,j)$.
We call $(i,j)$ the \emph{left neighbor} of $(i',j')$ if $i=i'$ and $j = j'-1$.

Let~$T$ be an SCT and $i,j\in T$ be two entries.
We refer to the \emph{row} and the \emph{column} of~$i$ in~$T$ by $r_T(i) \coloneqq r(T^{-1}(i))$ and $c_T(i)\coloneqq c(T^{-1}(i))$, respectively. 
We say that~$i$ \emph{attacks}~$j$ in~$T$ and write $i\att_T j$ if $T^{-1}(i) \att T^{-1}(j)$.
%The index~$T$ may be omitted if it is clear from context.
% We write $i\nei_T j$ if $T^{-1}(i)\nei T^{-1}(j).$ 
Note that $i\att_T j$ implies $i\neq j$. 
%Moreover,~$j$ is called \emph{left neighbor} of~$i$ if $T^{-1}(j)$ is the left neighbor of $T^{-1}(i)$.
If $T^{-1}(j)$ is the left neighbor of $T^{-1}(i)$ then we also call~$j$ the \emph{left neighbor} of~$i$.

For two sets of cells $C_1, C_2 \subseteq \N^2$ we say~$C_1$ attacks~$C_2$ and write $C_1 \att C_2$ if there are cells $c_1\in C_1$ and $c_2 \in C_2$ such that $c_1 \att c_2$. 
If $c(c_1)\leq c(c_2)$ for all $c_1\in C_1, c_2\in C_2$ then~$C_1$ is called \emph{left} of~$C_2$. If $c(c_1)< c(c_2)$ for all $c_1\in C_1, c_2\in c_2$,~$C_1$ is \emph{strictly left} of~$C_2$.
To simplify notation we may replace singletons by their respective element.
For instance, given a cell~$c_1$ we may write $c_1 \att C_2$ instead of $\set{c_1} \att C_2$.

In the same way we use these notions for sets of entries of an SCT.

\begin{exam}
		Consider the standard composition tableau~$T$ from~\cref{exa:sct}. We have $2\att_T 5, 3\att_T 4, 4\att_T 3$, $5\att_T 3$ and $i\natt_T j$ for all other pairs of entries. Moreover,~$3$ is left of $\set{1,4}$ in~$T$ and $2\att_T \set{3,5}$.
\end{exam}

\begin{defi}
	\label{thm:descents}
	 Let~$T$ be an SCT of size~$n$.
	\begin{enumerate}
		\item $\D(T) = \{i\in [n-1] \mid c_T(i) \leq c_T(i+1)\}$ is the \emph{descent set} of~$T$.
		\item $\AD(T) = \{i\in \D(T) \mid i\att_T i+1 \}$ is the \emph{set of attacking descents} of~$T$.
		\item $\nAD(T) = \{i\in \D(T) \mid i \notin \AD(T) \}$ is the \emph{set of non-attacking descents} of~$T$.
		\item[(1$'$)] $\DC(T) = \{i\in [n-1] \mid c_T(i+1) < c_T(i)\} = [n-1]\setminus \D(T)$ is the \emph{ascent set} of~$T$.
		\item[(2$'$)] $\ND(T) = \{i\in \DC(T) \mid i+1$ is the left neighbor of $i\}$ is the \emph{set of neighborly ascents} of~$T$.
		%\item $\nND(T) = \{i\in \DC(T) \mid i \notin \ND(T) \}$ is the \emph{set of non-neighborly ascents} of~$T$.
	\end{enumerate}
\end{defi}

\begin{exam}
Let~$T$ be the tableau from~\cref{exa:sct}. Then
	$\D(T) = \{2,3\}$, 
	$\AD(T) = \{3\}$, 
	$\DC(T) =\{1,4\}$ and
	$\ND(T) = \{4\}$. 
\end{exam}


\subsection{0-Hecke modules of standard composition tableaux}

In this subsection we consider the skew 0-Hecke modules $\sab$ and $\sse$ introduced in~\cite{Tewari2015} and review related results. This includes the special cases~$\sa$ and~$\se$.

\begin{thm}[{\cite[Theorem~9.8]{Tewari2015}}]
	\label{thm:sct_modules}
	Let~$\ab$ be a skew composition of size~$n$. 
	Then $\sab\coloneqq \spa_\K \SCT(\ab)$ is an~$\He$-module with respect to the following action. For $T\in \SCT(\ab)$ and $i=1,\dots, n-1$,
	\begin{align*}
	\pi_i T = 
	\begin{cases}
	T &\myif i \notin \D(T) \\
	0 &\myif i \in \AD(T) \\ 
	s_iT &\myif i \in \nAD(T) 
	\end{cases}
	\end{align*}
	where $s_iT$ is the tableau obtained from~$T$ by interchanging~$i$ and~$i+1$.
\end{thm}

The module~$\sa$ is called \emph{straight} if $\alpha= \ab$ is a composition.
Even though the main results of this paper are only for straight modules, we consider the more general concept of skew modules here as they naturally arise in the context of the 0-Hecke action on chains of $\mc L_c$ in~\cref{sec:0-hecke_operation_and_chains}.

\begin{exam}
	\label{exa:0-hecke_action}
	Consider the SCT 
	$T=
	\begin{ytableau} 
		1 \\ 
		6 & 5 & 4 & 3 \\ 
		8 & 7 & 2 \\ 
	\end{ytableau}.$ 
	Then $\D(T) = \set{1,2,6}$,	
	\[
		\pi_i T = 
		\begin{cases}
			T &\text{for $i=3,4,5,7$} \\
			0 &\text{for $i=6$} \\
			s_iT& \text{for $i=1,2$},
		\end{cases}
		\quad
		s_1T = 
		\begin{ytableau} 
			\bs2 \\ 
			6 & 5 & 4 & 3 \\ 
			8 & 7 & \bs1
		\end{ytableau}
		\quad \text{and} \quad
		s_2T =
		\begin{ytableau} 
			1 \\ 
			6 & 5 & 4 & \bs 2 \\ 
			8 & 7 & \bs3
		\end{ytableau}.
	\]
\end{exam}

We now decompose $\sab$ as in~\cite{Tewari2015}. To do this we use an equivalence relation.
Let~$\ab$ be a skew composition of size~$n$ and $T_1,T_2 \in \SCT(\ab)$. The equivalence relation $\sim$ on $\SCT(\ab)$ is given by
\begin{align*}
	T_1 \sim T_2 \iff 
	\text{in each column the relative orders of entries in~$T_1$ and~$T_2$ coincide.}
\end{align*}
 For example, the tableaux shown in~\cref{fig:example_for_straight_E} form an equivalence class under $\sim$ as well as the tableaux shown in~\cref{fig:example_for_skew_E}.
 We denote the \emph{set of equivalence classes} under $\sim$ on $\SCT(\ab)$ by $\mc E(\ab)$. 
 
For $E\in \mc E(\ab)$ define $\sse = \spa_\K E$. 
It is easy to see that the definition of the 0-Hecke action on standard composition tableaux in~\cref{thm:sct_modules} implies that $\sse$ is an~$\He$-submodule of $\sab$. Thus we have the following.
\begin{prop}[{\cite[Lemma~6.6]{Tewari2015}}]
	\label{thm:decomposition_in_equivalence_classes}
	Let~$\ab$ be a skew composition. Then we have
	$\sab = \bigoplus_{E\in \mc E(\ab)} \sse$ as $\He$-modules.
\end{prop}

The main result of this paper is that the $\He$-endomorphism ring of each straight module~$\se$ is $\K \id$ and, therefore, we obtain a decomposition of~$\sa$ into indecomposable submodules from~\cref{thm:decomposition_in_equivalence_classes}.


Let~$\ab$ be a skew composition of size~$n$ and $E\in \mc E(\ab)$.
 We continue by studying~$E$ and its module $\sse$ more deeply.
First, we consider a partial order $\pleq$ on~$E$.
It will turn out that $(E,\pleq)$ is a graded lattice.
Afterwards, we prepare two technical results,~\cref{thm:rank_in_E_and_column_word} and~\cref{thm:mutliple_flips_same_cell}, on the 0-Hecke action on standard composition tableaux for later use. 

Suppose $T_1,T_2\in E$. In~\cite[Section~4]{Tewari2015} it is shown that a partial order $\pleq$ on~$E$ is given by
\begin{align*}
	T_1 \pleq T_2 \iff \exists \sigma \in \fS_n \text{ such that } \pi_\sigma T_1 = T_2.
\end{align*}

 We refer to the poset $(E,\pleq)$ simply by~$E$. Two examples are shown in~\cref{fig:example_for_E}. The following theorem summarizes results of~\cite[Section~6]{Tewari2015}.


\begingroup 
\begin{figure}[htb]	\centering
%\def\captionseparator{}
%	\begin{subfigure}{.5\textwidth}
%\begin{minipage}{0.45\textwidth}\centering
%\def\figurename#1{}
%\def\thefigure{({\upshape\alph{figure}})}	\centering		
\subfigure[A poset $(E,\pleq)$ of straight tableaux]{%
		\scalebox{1}{%			
			\begin{tikzpicture} 
			\newcommand*{\xdist}{*3}
			\newcommand*{\ydist}{*2.2}
			
			\node (n0) at (0.00\xdist,0\ydist) 
			{
				$T_{0,E} =
				\begin{ytableau} 
				1 \\ 
				6 & 5 & 4 & 3 \\ 
				8 & 7 & 2 \\ 
				\end{ytableau} 
				$ 
			}; 
			\node (n1) at (-0.50\xdist,1\ydist) 
			{
				$\begin{ytableau} 
				2 \\ 
				6 & 5 & 4 & 3 \\ 
				8 & 7 & 1 \\ 
				\end{ytableau} 
				$ 
			}; 
			\node (n2) at (0.50\xdist,1\ydist) 
			{
				$\begin{ytableau} 
				1 \\ 
				6 & 5 & 4 & 2 \\ 
				8 & 7 & 3 \\ 
				\end{ytableau} 
				$ 
			}; 
			\node (n3) at (-0.50\xdist,2\ydist) 
			{
				$\begin{ytableau} 
				3 \\ 
				6 & 5 & 4 & 2 \\ 
				8 & 7 & 1 \\ 
				\end{ytableau} 
				$ 
			}; 
			\node (n4) at (0.50\xdist,2\ydist) 
			{
				$\begin{ytableau} 
				2 \\ 
				6 & 5 & 4 & 1 \\ 
				8 & 7 & 3 \\ 
				\end{ytableau} 
				$ 
			}; 
			\node (n5) at (-0.50\xdist,3\ydist) 
			{
				$\begin{ytableau} 
				4 \\ 
				6 & 5 & 3 & 2 \\ 
				8 & 7 & 1 \\ 
				\end{ytableau} 
				$ 
			}; 
			\node (n6) at (0.50\xdist,3\ydist) 
			{
				$\begin{ytableau} 
				3 \\ 
				6 & 5 & 4 & 1 \\ 
				8 & 7 & 2 \\ 
				\end{ytableau} 
				$ 
			}; 
			\node (n7) at (0.00\xdist,4\ydist) 
			{
				$T_{1,E} =
				\begin{ytableau} 
				4 \\ 
				6 & 5 & 3 & 1 \\ 
				8 & 7 & 2 \\ 
				\end{ytableau} 
				$ 
			}; 
			
			\draw [thick, ->] (n0) -- (n1) node [near start, left] {$\pi_{1}$}; 
			\draw [thick, ->] (n0) -- (n2) node [near start, left] {$\pi_{2}$}; 
			\draw [thick, ->] (n1) -- (n3) node [near start, left] {$\pi_{2}$}; 
			\draw [thick, ->] (n2) -- (n4) node [near start, left] {$\pi_{1}$}; 
			\draw [thick, ->] (n3) -- (n5) node [near start, left] {$\pi_{3}$}; 
			\draw [thick, ->] (n3) -- (n6) node [near start, left] {$\pi_{1}$}; 
			\draw [thick, ->] (n4) -- (n6) node [near start, left] {$\pi_{2}$}; 
			\draw [thick, ->] (n5) -- (n7) node [near start, left] {$\pi_{1}$}; 
			\draw [thick, ->] (n6) -- (n7) node [near start, left] {$\pi_{3}$}; 			
			\end{tikzpicture}				
		}\label{fig:example_for_straight_E}}
%	\caption{A poset $(E,\pleq)$ of straight tableaux}
	%	\end{subfigure}%
%\end{minipage}
%\begin{figure}
%	\begin{subfigure}{.5\textwidth}		
%\begin{minipage}{0.45\textwidth}
%\def\figurename{}\def\captionseparator{}
%\def\thefigure{({\upshape\alph{figure}})}	\centering		
\subfigure[A poset $(E', \pleq)$ of skew tableaux]{%
	\scalebox{1}{%	
		\newcommand*{\xdist}{*3}
		\newcommand*{\ydist}{*2.2}
		
		\begin{tikzpicture} 
		
		\node (n0) at (0.00\xdist,0.00\ydist) 
		{
			$T_{0,E'} = \begin{ytableau} 
			1 \\ 
			2 \\ 
			*(gray) & *(gray) & 4 & 3 \\ 
			\end{ytableau} 
			$ 
		}; 
		\node (n1) at (0.00\xdist,1.00\ydist) 
		{
			$\begin{ytableau} 
			1 \\ 
			3 \\ 
			*(gray) & *(gray) & 4 & 2 \\ 
			\end{ytableau} 
			$ 
		}; 
		\node (n2) at (-0.50\xdist,2.00\ydist) 
		{
			$\begin{ytableau} 
			1 \\ 
			4 \\ 
			*(gray) & *(gray) & 3 & 2 \\ 
			\end{ytableau} 
			$ 
		}; 
		\node (n3) at (0.50\xdist,2.00\ydist) 
		{
			$\begin{ytableau} 
			2 \\ 
			3 \\ 
			*(gray) & *(gray) & 4 & 1 \\ 
			\end{ytableau} 
			$ 
		}; 
		\node (n4) at (0.00\xdist,3.00\ydist) 
		{
			$\begin{ytableau} 
			2 \\ 
			4 \\ 
			*(gray) & *(gray) & 3 & 1 \\ 
			\end{ytableau} 
			$ 
		}; 
		\node (n5) at (0.00\xdist,4.00\ydist) 
		{
			$T_{1,E'} = \begin{ytableau} 
			3 \\ 
			4 \\ 
			*(gray) & *(gray) & 2 & 1 \\ 
			\end{ytableau} 
			$ 
		}; 
		
		\draw [thick, ->] (n0) -- (n1) node [near start, left] {$\pi_{2}$}; 
		\draw [thick, ->] (n1) -- (n2) node [near start, left] {$\pi_{3}$}; 
		\draw [thick, ->] (n1) -- (n3) node [near start, left] {$\pi_{1}$}; 
		\draw [thick, ->] (n2) -- (n4) node [near start, left] {$\pi_{1}$}; 
		\draw [thick, ->] (n3) -- (n4) node [near start, left] {$\pi_{3}$}; 
		\draw [thick, ->] (n4) -- (n5) node [near start, left] {$\pi_{2}$}; 
		
		\end{tikzpicture} 

	}		\label{fig:example_for_skew_E}
}
%		\caption{A poset $(E', \pleq)$ of skew tableaux}
	%		\end{minipage}
%	\end{subfigure}
%\def\captionseparator{.}
%\setcounter{figure}{0}
		\caption{%
			Two posets each of them given by an equivalence class of standard composition tableaux and the corresponding partial order $\pleq$.
			Each covering relation is labeled with the 0-Hecke generator $\pi_i$ realizing it.
		}
		\label{fig:example_for_E}
%\vspace*{-2mm}
\end{figure}
\endgroup


\begin{theo}
	\label{thm:source_and_sink_tableau}
	 Let~$\ab$ be a skew composition, $E\in \mc E(\ab)$ and $T\in E$.
	\begin{enumerate}
		\item The tableau $T$ is minimal according to $\pleq$ if and only if $\DC(T)= \ND(T)$. There is a unique tableau $T_{0,E}\in E$ which satisfies these conditions called \emph{source tableau} of~$E$.
		\item The tableau $T$ is maximal according to $\pleq$ if and only if $\D(T)= \AD(T)$. There is a unique tableau $T_{1,E}\in E$ which satisfies these conditions called \emph{sink tableau} of~$E$.
	\end{enumerate}
	In particular, $\sse$ is a cyclic module generated by $T_{0,E}$.
\end{theo}

Source and sink tableaux can be observed in~\cref{fig:example_for_E}. Next, we establish a connection between~$E$ and an interval of the left weak order. To do this we use the notion of \emph{column words}.
 Given $T\in \SCT(\ab)$ and $j\geq 1$, let~$w_j$ be the word obtained by reading the entries in the~$j$th column of~$T$ from top to bottom. Then $\col_T= w_1 w_2 \cdots$ is the \emph{column word} of~$T$. Clearly, $\col_T$ can be regarded as an element of $\fS_n$ (in one-line notation).
 
\begin{exam}
	The column word of the tableau $T_{0,E}$ from~\cref{fig:example_for_straight_E} is given by $\col_{T_{0,E}} = 16857423\in \fS_8$.
\end{exam}

\begin{lemma}[{\cite[Lemma~4.4]{Tewari2015}}]
	\label{thm:tableaux_cover_impies_column_words_cover}
	Let~$T_1$ be an SCT, $i\in \nAD(T_1)$ and $T_2 = \pi_i T_1$. Then $\col_{T_2} = \fs_i \col_{T_1}$ and $\ell(\col_{T_2}) = \ell(\col_{T_1}) + 1$. That is, $\col_{T_2}$ covers $\col_{T_1}$ in left weak order.
\end{lemma}


The following statement is similar to~\cite[Lemma~4.3]{Tewari2015}.
\begin{lemma}
	\label{thm:column_word_and_0-hecke_operation}
\looseness-1
	Let~$T_1$ and~$T_2$ be two standard composition tableaux such that $\pi_{i_p} \cdots \pi_{i_1} T_1 = T_2$. Then there is a subsequence $j_q, \ldots, j_1$ of $i_p, \ldots,i_1$ such that
	\begin{enumerate}
		\item
		$T_2 = \pi_{j_q} \cdots \pi_{j_1}T_1$,
		\item 
		$\fs_{j_q} \cdots \fs_{j_1}$ is a reduced word for $\col_{T_2}\col^{-1}_{T_1}$.
	\end{enumerate}
	In particular, $T_2 = \pi_{\col_{T_2}\col^{-1}_{T_1}} T_1$.
\end{lemma}


\begin{proof} It follows from the definition of the 0-Hecke operation that we can find a subsequence $j_q, \ldots, j_1$ of $i_p, \ldots,i_1$ of minimal length such that $T_2 = \pi_{j_q} \cdots \pi_{j_1}T_1$. If~$q=0$ then $T_2 = T_1$ and the result is trivial. If $q=1$ set $i \coloneqq j_1$. Then by the minimality of~$q$, $T_2 \neq T_1$ and thus $i\in \nAD(T_1)$. Now~\cref{thm:tableaux_cover_impies_column_words_cover} shows that $\fs_i$ is a reduced word for $\col_{T_2}\col^{-1}_{T_1}$. If $q>1$ use the case $q=1$ iteratively.
\end{proof}



\begin{thm}[{\cite[Theorem~6.18]{Tewari2015}}]
	\label{thm:E_isomorphic_to_bruhat_interval}
	Let $\ab$ be a skew composition, $E\in \mc E(\ab)$ and $I = [\col_{T_{0,E}}, \col_{T_{1,E}}]$ an interval in left weak order. Then the map 
	$\col \colon E\to I$, $T\mapsto\col_T$ is a poset isomorphism. In particular,~$E$ is a graded lattice with rank function $\delta\colon T\mapsto \ell(\col_T \col^{-1}_{T_{0,E}})$.
\end{thm}

Actually,~\cref{thm:source_and_sink_tableau},~\cref{thm:tableaux_cover_impies_column_words_cover} and~\cref{thm:column_word_and_0-hecke_operation} are everything needed to prove~\cref{thm:E_isomorphic_to_bruhat_interval} as in~\cite{Tewari2015}. They imply that $\col$ (and its inverse) map maximal chains to maximal chains.
Note that it follows from~\cref{thm:E_isomorphic_to_bruhat_interval} and~\cref{thm:properties_bruhat_order} that for $T_1 \pleq T_2$ saturated chains from~$T_1$ to~$T_2$ correspond to reduced words for $\col_{T_2}\col^{-1}_{T_1}$.


\begin{coro}
	\label{thm:rank_in_E_and_column_word}
	Let~$T_1$ and~$T_2$ be two standard composition tableaux of size~$n$ and $\sigma \in \fS_n$ such that $T_2 = \pi_\sigma T_1$. Then~$T_1$ and~$T_2$ belong to the same equivalence class under $\sim$. Let $\delta$ be the rank function of that class. Then
	\begin{enumerate}
		\item\label{coro2.24_1} $\delta(T_2) - \delta(T_1) = \ell(\col_{T_2}\col^{-1}_{T_1})$,
		\item\label{coro2.24_2} $\delta(T_2) - \delta(T_1) \leq \ell(\sigma)$ where we have equality if and only if $\sigma = \col_{T_2}\col_{T_1}^{-1}$.
	\end{enumerate}
\end{coro}


\begin{proof} Since $T_2 = \pi_\sigma T_1$, $T_2 \sim T_1$.
We obtain part~\eqref{coro2.24_1} from the above discussion. Part~\eqref{coro2.24_2} is a consequence of~\eqref{coro2.24_1} and~\cref{thm:column_word_and_0-hecke_operation}.
\end{proof}

\looseness-1
%We finish this section by preparing another consequence of~\cref{thm:tableaux_cover_impies_column_words_cover} for
We finish this section by preparing another result for~\cref{sec:endomorphism_ring_of_0-Hecke_modules}.

\begin{prop}
	\label{thm:mutliple_flips_same_cell}
	
	Let~$T$ be an SCT, $i,j\in T$ be such that $i<j$ and $\cell = T^{-1}(i)$. If in~$T$~$i$ is located left of $[i+1,j]$
	and does not attack $[i+1,j]$ then 
	\begin{enumerate}	
		\item\label{prop2.25_1} $T'\coloneqq \pi_{j-1} \cdots \pi_{i+1}\pi_i T$ is an SCT,
		\item\label{prop2.25_2} $\fs_{j-1} \cdots \fs_{i+1}\fs_i$ is a reduced word for $\col_{T'}\col_T^{-1}$,
		\item\label{prop2.25_3} $T'(\cell)= j$.
	\end{enumerate}
\end{prop}

\begin{proof}
	Let~$T$ be an SCT and $i,j\in T$ such that $i<j$,~$i$ is located left of $[i+1,j]$
	and $i \natt[i+1,j]$. Set $\cell = T^{-1}(i)$.
	We do an induction on $m\coloneqq j-i$. If $m=1$ then $i\in \nAD(T)$ and $T' = \pi_i T$. Thus, 	\eqref{prop2.25_1} and~\eqref{prop2.25_3} hold by the definition of the 0-Hecke action and~\eqref{prop2.25_2} is a consequence of~\cref{thm:tableaux_cover_impies_column_words_cover}.
	
	Now, let $m>1$. Since by assumption~$i$ is located left of $[i+1,j]$
	and $i~\natt~[i+1,j]$, we can apply the induction hypothesis on~$i$ and $j-1$ and obtain that $T''\coloneqq \pi_{j-2} \cdots \pi_{i+1}\pi_i T$ is an SCT, $\fs_{j-2} \cdots \fs_{i+1}\fs_i$ is a reduced word for $\col_{T''}\col_T^{-1}$ and $T''(\cell) = j-1$.
	Since the operators $\pi_{j-2}, \ldots, \pi_{i+1}, \pi_{i}$ are unable to move~$j$, we have $T''^{-1}(j) = T^{-1}(j).$
	By choice of~$i$ and~$j$, $\cell \natt T^{-1}(j) = T''^{-1}(j)$ and $\cell$ is left of $T''^{-1}(j)$.
	Thus, $j-1\in \nAD(T'')$ so that $T'~=~\pi_{j-1} \pi_{j-2}\cdots \pi_i T~=~\pi_{j-1}(T'')$ is an SCT and $T'(\cell) = j$.
	It follows from~\cref{thm:tableaux_cover_impies_column_words_cover} that $\col_{T'}\col_T^{-1}=s_{j-1}\col_{T''}\col^{-1}_{T}= \fs_{j-1}\fs_{j-2} \cdots \fs_i$ and that $\fs_{j-1}\fs_{j-2} \cdots \fs_i$ is a reduced word.
\end{proof}


